\documentclass[12pt]{article}
\usepackage[utf8]{inputenc}
\usepackage[T1]{fontenc}
\usepackage{hyperref}
\title{Nhengatu Language Resources for K-12 Curriculum Development}
\author{Compiled by Claude Sonnet 5.5 (claude-sonnet-5-5) for Andr\'as Kornai}
\date{}
\begin{document}
\maketitle

\section{Introduction}
Nhengatu, also written Nheengatu (ISO 639-3 yrl; Glottocode nhen1239), is a Tupian language of the upper Rio Negro region of Brazil, with speakers also in Colombia and Venezuela. Speaker estimates in the search results diverge widely: about 6,000, 8,000 and 14,000. It was co-official, with Tukano and Baniwa, in S\~ao Gabriel da Cachoeira (Amazonas) under Municipal Law 145 of 2002; Yanomami was added in 2017.

\section{Primary Educational Resources}
\begin{itemize}
\item UFSCar (Leetra) publications include bilingual materials, among them a medicinal plants booklet in Nheengatu and a 2014 reading and writing book in Y\~eg\'at\'u made with pupils of the Escola Municipal Ind\'igena ``Pastor Jaime''. The search did not show which of the two journal issues holds which item: \url{https://www.leetra.ufscar.br/files/biblioteca/revista_li_15.pdf} and \url{https://www.leetra.ufscar.br/files/biblioteca/revista_li_especial_3.pdf}
\item IFAM Language Center, 180-hour Nheengatu course in the Amazonas municipality: \url{https://portalamazonia.com/educacao/nheengatu-e-ofertado-por-centro-de-idiomas-do-ifam-em-municipio-do-amazonas/}
\end{itemize}
A UFAM dissertation reports that teaching of the co-official languages in schools remains limited in form and in quantity, with material obstacles to implementation: \url{https://tede.ufam.edu.br/handle/tede/11209}

\section{Linguistic and Reference Materials}
\begin{itemize}
\item Eduardo de Almeida Navarro, \emph{Curso de L\'ingua Geral (Nheengatu ou Tupi Moderno)}, 2011, a course and grammar aimed at revitalisation (USP): \url{https://tupi.fflch.usp.br/sites/tupi.fflch.usp.br/files/CURSO%20DE%20L\%C3\%8DNGUA%20GERAL%20(NHEENGATU).pdf}
\item Yam\~a, Yaguak\~ag, Reis and Jos\'e (2018), \emph{Dicion\'ario Nh\~e\~egat\'u Tradicional} (Nova Olinda do Norte, AM). I saw it only cited in a search summary and found no link to it.
\end{itemize}

\section{Religious and Community Resources}
None found in the search results.

\section{Digital and Audio Resources}
\begin{itemize}
\item Translation app and platform for Nheengatu built by researchers (UFC, Complin group): \url{https://diariodonordeste.verdesmares.com.br/ceara/nheengatu-cearenses-criam-plataforma-para-traducao-da-unica-lingua-viva-descendente-do-tupi-antigo-1.3229012}
\item Press report on the same app (IHU Unisinos, 2022): \url{https://ihu.unisinos.br/categorias/621911-pesquisadores-criam-aplicativo-para-traducao-da-unica-lingua-viva-descendente-do-tupi}
\end{itemize}

\section{Curriculum Development Context}
Official status in S\~ao Gabriel da Cachoeira gives a legal basis for school use, but sources describe implementation as limited. Bilingual teaching material exists in small quantity, and the grammar by Navarro and the dictionary of 2018 can support materials development. The most useful next step would be to contact the municipal education secretariat and the indigenous schools of the Rio Negro. Sources disagree on speaker numbers, and no Nheengatu-specific orthography policy was found.

\end{document}